% sections/02-escaneo-nmap.tex
\section{Fase de escaneo}

\subsection{Identificación del objetivo}

Para la fase de escaneo se hizo uso de la máquina virtual Metasploitable, cuya dirección IP se identificó desde la propia máquina con el comando \texttt{ifconfig} (la interfaz \texttt{eth0} entrega la dirección de la red local sobre la que se ejecutan todos los escaneos de esta sección).

\porque{IP de Metasploitable}{%
  Anota en el texto la IP real de su Metasploitable y reemplaza
  \texttt{192.168.1.116} en todos los comandos y capturas de las secciones 2
  a 5 (y \texttt{192.168.1.118} por la IP de su Kali en la sección 4); la
  guía advierte que las IPs pueden cambiar entre montajes. Si quieren
  evidenciar este paso, añadan una captura adicional
  (\texttt{pics/metasploitable-ifconfig.png}) con la salida del ifconfig.%
}

\subsection{Escaneo sencillo}

Desde la máquina de Kali Linux se ejecutó un escaneo sencillo sobre el objetivo, reemplazando la IP mostrada por la de la máquina Metasploitable usada:

\begin{lstlisting}[language=bash]
nmap 192.168.1.116
\end{lstlisting}

\captura{nmap-sencillo.png}{Escaneo sencillo de Nmap sobre Metasploitable: puertos TCP abiertos y servicios asociados}

\begin{analisis}
\end{analisis}

\porque{qué debe verse en la captura y el análisis}{%
  El escaneo por defecto (TCP connect hacia los mil puertos más comunes)
  lista los puertos abiertos con su estado y el servicio asociado; para
  Metasploitable se espera una lista larga (21/ftp, 22/ssh, 23/telnet,
  25/smtp, 80/http, 139/445/samba, entre otros, con \texttt{mysql},
  \texttt{postgresql}, \texttt{vnc} o \texttt{distccd} según el rango
  escaneado). Base para el análisis: ``La salida evidencia un servidor con
  una cantidad anormal de servicios expuestos, varios de ellos de
  administraci\'on remota (telnet, ftp, vnc), lo que delata un entorno de
  pruebas deliberadamente vulnerable; cada puerto abierto amplía la
  superficie de ataque disponible para la siguiente fase.'' Verifica contra
  tu salida qué puertos alcanzó a reportar el escaneo por defecto (solo son
  los mil más conocidos) y coméntalo.%
}

\subsection{Métodos de escaneo}

\pregunta{Complete el siguiente cuadro describiendo brevemente en qué consisten los métodos de escaneo de Nmap.}
\begin{respuesta}
  \begin{center}
  \small
  \begin{tabularx}{\linewidth}{>{\ttfamily}l X}
    \toprule
    \normalfont\bfseries Opción & \bfseries Definición \\
    \midrule
    -sS & Sondeo TCP SYN (\eng{half-open}): se envía un paquete SYN a cada puerto y, de recibir SYN/ACK (puerto abierto), se responde con RST abortando el handshake antes de completarlo, de forma que muchos servicios no llegan a registrarlo; es rápido y relativamente discreto. \\
    -sT & Sondeo TCP connect(): establece la conexión completa mediante la llamada al sistema \texttt{connect()}, con lo cual funciona sin privilegios especiales y es compatible con cualquier sistema, a costa de ser más lento y dejar rastro en los registros de auditoría de cada servicio. \\
    -sU & Sondeo UDP: envía un datagrama (vacío por defecto) a cada puerto y clasifica según la respuesta (respuesta UDP = abierto, ICMP puerto inalcanzable = cerrado, sin respuesta = abierto o filtrado); es el más lento por la naturaleza del protocolo y la limitación del ICMP. \\
    -sN / -sF / -sX & Sondeos TCP Null, FIN y Xmas: aprovechando el RFC 793, envían paquetes sin bandera alguna (Null), solo FIN, o FIN+PSH+URG (\eng{Xmas}), de forma que un puerto cerrado responde RST mientras que uno abierto ignora el paquete; al no establecerse conexión no generan registros, aunque los cortafuegos con estado los detectan y filtran. \\
    -sA & Sondeo TCP ACK: envía paquetes con la bandera ACK (como si respondieran a una conexión existente) para sondear las reglas del cortafuegos; un RST indica el estado sin filtrar, mientras que la ausencia de respuesta o un ICMP inalcanzable delata el filtrado, sin distinguir puertos abiertos de cerrados. \\
    -sW & Sondeo de ventana TCP: variante del sondeo ACK que examina el tamaño de la ventana TCP reportado en el RST recibido, pues ciertos sistemas devuelven ventanas positivas para los puertos abiertos y nulas para los cerrados, con lo que en esos casos llega a diferenciar el estado. \\
    -sM & Sondeo TCP Maimon: envía paquetes FIN/ACK a cada puerto; muchos sistemas derivados de BSD devuelven RST para los puertos abiertos, de forma que en esos sistemas permite distinguirlos de los cerrados, compartiendo con Null/FIN/Xmas la ventaja de no generar conexiones registrables. \\
    \bottomrule
  \end{tabularx}
  \end{center}

  Los métodos marcados con asterisco en la guía (sS, sU, sN/sF/sX, sA, sW y
  sM) requieren privilegios de root, pues construyen paquetes crudos
  (\eng{raw packets}) en lugar de usar la pila TCP del sistema.
\end{respuesta}

\porque{antes de continuar con la tabla}{%
  Definiciones tomadas de la guía de referencia de Nmap \cite{nmap-book}:
  revíalas contra el \texttt{man nmap} de su Kali y quita o matiza lo que no
  coincida con lo visto en clase. Ojo con el encoding: copia la tabla tal
  cual, ya está escrita para no romper el ancho de columna.%
}

Sobre el mismo objetivo se ejecutaron escaneos con los métodos descritos, adjuntando los pantallazos y sus resultados:

\captura{nmap-metodos.png}{Escaneos ejecutados con los distintos métodos de Nmap sobre Metasploitable y resultados obtenidos}

\begin{analisis}
\end{analisis}

\porque{evidencia de los escaneos}{%
  La guía pide pantallazos de los escaneos ejecutados: si quieren evidenciar
  cada método por separado, añadan una captura por ejecución con nombres
  tipo \texttt{nmap-sS.png}, \texttt{nmap-sU.png}, \texttt{nmap-sX.png},
  etc. (la sintaxis es \texttt{nmap -sS 192.168.1.116}).   Para el análisis
  compara: los métodos SYN/NULL/FIN/Xmas requieren root y devuelven los
  mismos abiertos, el ACK clasifica casi todo como sin filtrar (está mapeando
  el cortafuegos, no los puertos) y el UDP es mucho más lento y deja casi
  todo en abierto|filtrado al no haber respuesta.%
}

\pregunta{¿De qué forma podemos usar NMAP para evitar un Firewall (cortafuegos)? Describa brevemente las opciones que podemos configurar en la herramienta y cómo funcionaría.}
\begin{respuesta}
\end{respuesta}

\porque{opciones con las que armar la respuesta}{%
  Lista con término en negrita + explicación, tu patrón usual, apoyada en
  \cite{nmap-book}: \textbf{Fragmentación (-f):} parte el paquete sondeador
  en fragmentos diminutos para que un cortafuegos simple no logre
  reensamblar y comparar el paquete completo; \textbf{Señuelos (-D
  RND:10):} intercala sondeos con direcciones IP falsas junto a la real,
  de forma que los registros del objetivo no revelan cuál es el origen
  verdadero; \textbf{Suplantación de origen (-S) y de MAC (--spoof-mac):}
  usan una dirección distinta para los paquetes (útil en redes filtradas por
  dirección o por fabricante); \textbf{Puerto de origen (--source-port 53):}
  fija el puerto de origen de los sondeos en uno de confianza (DNS, 53),
  pues muchos cortafuegos permiten el tráfico de regreso a servicios
  conocidos; \textbf{Temporización (-T0/T1, --max-rtt-timeout):} espacia los
  sondeos para no disparar la detección por tasa de paquetes de un
  IDS/IPS; \textbf{Sondeo SYN (-sS) en lugar de connect():} al no completar
  la conexión, los servicios no lo registran en sus auditorías; cierra con
  la idea de que se trata de evadir \emph{detección}, no de atravesar
  cualquier cortafuegos (los con estado modernos hacen mucho menos eficaz
  la técnica).%
}

\pregunta{¿Qué importancia puede tener el reconocimiento de puertos dentro de la fase de Escaneo?}
\begin{respuesta}
\end{respuesta}

\porque{base para la respuesta}{%
  El puerto abierto es la puerta concreta a la que un atacante puede
  llamar: identificado el puerto y su servicio (y mejor aún su versión,
  con \texttt{-sV}), se puede buscar el exploit específico de esa versión,
  dirigir fuerza bruta contra servicios de autenticación o detectar
  servicios de administración mal expuestos; en términos de la metodología,
  el escaneo de puertos convierte la lista de hosts vivos del
  reconocimiento en una superficie de ataque explotable en la fase
  siguiente.%
}

\pregunta{Poniéndose en la situación de un atacante, ¿qué tipo de ataques considera que pueden llegarse a realizarse con la información levantada hasta el momento?}
\begin{respuesta}
\end{respuesta}

\porque{base para la respuesta}{%
  Con los servicios y versiones identificadas: explotación de
  vulnerabilidades conocidas de cada versión (la que precisamente se hace
  en las secciones 4 y 5 con distcc y vsftpd), fuerza bruta de credenciales
  contra ssh/telnet/ftp/mysql, accesos con credenciales por defecto,
  enumeración de usuarios y recursos compartidos (Samba, NFS) y uso del
  servidor como punto de apoyo para movimientos laterales dentro de la red;
  conecta la respuesta con lo que el reporte de la sección 3 confirma
  después.%
}
